\section{Randomized --- UPath Machine}
%%% generated by the blueprint pipeline from TCSlib.Complexity.Randomized.UPathMachine
\subsection*{Randomized --- UPath Machine: Introduction}

Declarations of the Lean module \texttt{TCSlib.Complexity.Randomized.UPathMachine}.

\subsection*{Randomized --- UPath Machine: Results}

\begin{definition}[Compression program sites]
\lean{Complexity.LogProg.CpSite}
\label{Complexity.LogProg.CpSite}
\leanok
Introduce a finite set of eleven program-site labels: $yq$, $sTmp$, $sP$, $sCur$, $tQ$,
$tTmp$, $tP$, $tT$, $wTmp$, $wNx$, and $wCur$.
\end{definition}

\begin{definition}[Index sites]
\lean{Complexity.LogProg.IxSite}
\label{Complexity.LogProg.IxSite}
\leanok
The index-site set consists of four distinct sites: $y_0$, $y_1$, $y_2$, and $w$.
\end{definition}

\begin{definition}[Addition sites]
\lean{Complexity.LogProg.AdSite}
\label{Complexity.LogProg.AdSite}
\leanok
\uses{Complexity.LogProg.IxSite}
The addition-site set consists of sites indexed by an index site and one of three
positions, five further sites indexed by one of five positions, and one additional site.
\end{definition}

\begin{definition}[Multiplication sites]
\lean{Complexity.LogProg.MulSite}
\label{Complexity.LogProg.MulSite}
\leanok
The multiplication-site set consists of two distinct sites, $m_0$ and $m_1$.
\end{definition}

\begin{definition}[Machine labels]
\lean{Complexity.LogProg.ULabel}
\label{Complexity.LogProg.ULabel}
\leanok
\uses{Complexity.LogProg.AdSite, Complexity.LogProg.CpSite, Complexity.LogProg.IxSite, Complexity.LogProg.MulSite}
The label set consists of the listed named labels, together with labels indexed by
positions in a five-element set, by multiplication sites, by copy sites, by addition
sites, or by index sites. For each addition site there are four indexed labels, and for
each index site there are nine indexed labels.
\end{definition}

\begin{definition}[Copy-site destinations]
\lean{Complexity.LogProg.cpDst}
\label{Complexity.LogProg.cpDst}
\leanok
\uses{Complexity.LogProg.CpSite}
Assign to the copy sites $yq,sTmp,sP,sCur,tQ,tTmp,tP,tT,wTmp,wNx,wCur$ the respective
destination positions $9,6,0,4,9,6,0,5,6,8,4$ in a fifteen-position row.
\end{definition}

\begin{definition}[Copy-site sources]
\lean{Complexity.LogProg.cpSrc}
\label{Complexity.LogProg.cpSrc}
\leanok
\uses{Complexity.LogProg.CpSite}
Assign to the copy sites $yq,sTmp,sP,sCur,tQ,tTmp,tP,tT,wTmp,wNx,wCur$ the respective
source positions $0,11,9,11,0,11,9,11,1,13,8$ in a fifteen-position row.
\end{definition}

\begin{definition}[Copy-site label assignment]
\lean{Complexity.LogProg.cpEx}
\label{Complexity.LogProg.cpEx}
\leanok
\uses{Complexity.LogProg.CpSite, Complexity.LogProg.ULabel}
Let the copy sites be $yq$, $sTmp$, $sP$, $sCur$, $tQ$, $tTmp$, $tP$, $tT$, $wTmp$,
$wNx$, and $wCur$. Assign them, respectively, the labels $y0$, $s3$, $stt$, $t0$, $t1$,
$t4$, $ttt$, $lc(0)$, $c0$, $c4$, and $stay$.
\end{definition}

\begin{definition}[Row positions of index sites]
\lean{Complexity.LogProg.ixRow}
\label{Complexity.LogProg.ixRow}
\leanok
\uses{Complexity.LogProg.IxSite}
Let the index-site set be $\{y_0,y_1,y_2,w\}$. Define a position in
$\{0,1,\ldots,14\}$ by assigning $y_0$ to $11$, $y_1$ and $y_2$ to $12$, and $w$ to
$4$.
\end{definition}

\begin{definition}[Column assignment to index sites]
\lean{Complexity.LogProg.ixCol}
\label{Complexity.LogProg.ixCol}
\leanok
\uses{Complexity.LogProg.IxSite}
Let the four index sites be $y_0,y_1,y_2,w$. Map them to column indices in
$\{0,\ldots,14\}$ by sending $y_0$ to $12$, $y_1$ and $y_2$ to $11$, and $w$ to
$8$.
\end{definition}

\begin{definition}[Index site label assignment]
\lean{Complexity.LogProg.ixEx}
\label{Complexity.LogProg.ixEx}
\leanok
\uses{Complexity.LogProg.IxSite, Complexity.LogProg.ULabel}
Define a function from the four index sites to labels by sending the sites $y_0$, $y_1$,
and $y_2$ to the labels $y_4$, $y_{6t}$, and $y_{6f}$, respectively, and sending the
site $w$ to the word label.
\end{definition}

\begin{definition}[Addition site columns]
\lean{Complexity.LogProg.adY}
\label{Complexity.LogProg.adY}
\leanok
\uses{Complexity.LogProg.AdSite, Complexity.LogProg.ixCol}
Define a map from addition sites to column indices $0$ through $14$ as follows: an index
site $(k,j)$ maps to $1$ when $j=0$ or $1$, and to the column assigned to index site $k$
when $j=2$; every site $(	ext{ten},j)$ maps to $1$, and the doubling site maps to $8$.
\end{definition}

\begin{definition}[Addition site positions]
\lean{Complexity.LogProg.adP}
\label{Complexity.LogProg.adP}
\leanok
\uses{Complexity.LogProg.AdSite}
Define a map from addition sites to column indices $0$ through $14$ by sending every
index site $(k,j)$ to $0$, every site $(	ext{ten},j)$ to $13$, and the doubling site to
$13$.
\end{definition}

\begin{definition}[Addition-site label assignment]
\lean{Complexity.LogProg.adEx}
\label{Complexity.LogProg.adEx}
\leanok
\uses{Complexity.LogProg.AdSite, Complexity.LogProg.ULabel, Complexity.LogProg.ixEx}
Define a map from addition sites to labels as follows. For an index site $(k,j)$, assign
the label indexed by $k$ with subscript $2$ when $j=0$, with subscript $8$ when $j=1$,
and the example label assigned to $k$ when $j=2$. The five sites indexed by
$j=0,1,2,3,4$ map respectively to $	ext{lc}_1$, $	ext{lc}_2$, $	ext{lc}_3$,
$	ext{lc}_4$, and $	ext{mc}_{m_0}$; the doubling site maps to $c_3$.
\end{definition}

\begin{definition}[Multiplication site indices]
\lean{Complexity.LogProg.mulA}
\label{Complexity.LogProg.mulA}
\leanok
\uses{Complexity.LogProg.MulSite}
Let the two multiplication sites be $m_0$ and $m_1$. Define a map from these sites to
the indices $0,1,\ldots,14$ by sending $m_0$ to $1$ and $m_1$ to $14$.
\end{definition}

\begin{definition}[Multiplication site indices]
\lean{Complexity.LogProg.mulB}
\label{Complexity.LogProg.mulB}
\leanok
\uses{Complexity.LogProg.MulSite}
Let the two multiplication sites be $m_0$ and $m_1$. Define their indices in
$	ext{Fin}(15)$ by assigning $13$ to $m_0$ and $14$ to $m_1$.
\end{definition}

\begin{definition}[Right column indices for multiplication sites]
\lean{Complexity.LogProg.mulR}
\label{Complexity.LogProg.mulR}
\leanok
\uses{Complexity.LogProg.MulSite}
Let $M$ be the two-element set of multiplication sites, with elements $m_0$ and $m_1$.
Define a map $r:M\to\{0,1,\ldots,14\}$ by $r(m_0)=14$ and $r(m_1)=10$.
\end{definition}

\begin{definition}[Successor labels for multiplication sites]
\lean{Complexity.LogProg.mulEx}
\label{Complexity.LogProg.mulEx}
\leanok
\uses{Complexity.LogProg.MulSite, Complexity.LogProg.ULabel}
Define a map from the two multiplication sites to labels by sending $m_0$ to the label
$mc(m_1)$ and $m_1$ to the label $w_0$.
\end{definition}

\begin{definition}[Universal path register program]
\lean{Complexity.LogProg.upA}
\statementsource{arora-barak-ch07-randomized-computation}{pdf-26d95080dbde-p016-b002}
\label{Complexity.LogProg.upA}
\leanok
\uses{Complexity.LogProg.NARM, Complexity.LogProg.ULabel, Complexity.LogProg.adEx, Complexity.LogProg.adP, Complexity.LogProg.adY, Complexity.LogProg.cpDst, Complexity.LogProg.cpEx, Complexity.LogProg.cpSrc, Complexity.LogProg.ixRow, Complexity.LogProg.mulA, Complexity.LogProg.mulB, Complexity.LogProg.mulEx, Complexity.LogProg.mulR}
Define a nondeterministic register program with fifteen registers, two deciders, and the
specified control labels, whose instruction at each label is the one given by the
program's instruction assignment.
\end{definition}

\begin{lemma}[Safety Invariant for the Program]
\lean{Complexity.LogProg.upA_npreS}
\label{Complexity.LogProg.upA_npreS}
\leanok
\uses{Complexity.LogProg.AConf, Complexity.LogProg.NARM.detPart, Complexity.LogProg.NPreS, Complexity.LogProg.PreS, Complexity.LogProg.ULabel, Complexity.LogProg.upA}
\difficulty{5}
For every Boolean input $x$ and abstract configuration consisting of a control state,
fifteen natural-number registers, and an optional Boolean result, the following
condition holds. If the state is live and its instruction compares two registers, their
indices are distinct; if it compares a register with the first payload of a plain
encoded input or with either payload of a nested encoded input, respectively, then the
input has the corresponding valid canonical encoding; and if it calls a decider, its
argument-register list has no repetitions. No further condition is imposed in all other
cases, including halted configurations and nondeterministic choice states.
\end{lemma}

\begin{definition}[Register states for the logarithmic program]
\lean{Complexity.LogProg.UReg}
\label{Complexity.LogProg.UReg}
\leanok
A register state is a tuple of fifteen natural numbers, denoted
$p,n,a,b,cur,t,tmp,st,nx,q,L,u,w,e,f$, representing respectively the input pointer,
vertex count, outer counter, inner counter, current vertex, target, scratch copy, step
count, random index, saved field position, walk length, candidate vertex, column,
scratch sum, and scratch product.
\end{definition}

\begin{definition}[Register-value vector]
\lean{Complexity.LogProg.UReg.v}
\label{Complexity.LogProg.UReg.v}
\leanok
\uses{Complexity.LogProg.UReg}
For a register state, define its value vector on the indices $\{0,\ldots,14\}$ by listing, in order,
its fifteen fields $p,n,a,b,cur,t,tmp,st,nx,q,L,u,w,e,f$.
\end{definition}

\begin{definition}[Register bound]
\lean{Complexity.LogProg.UReg.Le}
\label{Complexity.LogProg.UReg.Le}
\leanok
\uses{Complexity.LogProg.UReg}
For a register state and a natural-number bound $B$, the state is bounded by $B$ exactly
when each of its fifteen fields is at most $B$.
\end{definition}

\begin{lemma}[Register Bound for Abstract Configurations]
\lean{Complexity.LogProg.UReg.reg}
\label{Complexity.LogProg.UReg.reg}
\leanok
\uses{Complexity.LogProg.RegLe, Complexity.LogProg.ULabel, Complexity.LogProg.UReg, Complexity.LogProg.UReg.Le, Complexity.LogProg.UReg.v}
\difficulty{4}
Let $\sigma$ be a register state, let $B$ be a natural number, let $l$ be an optional
program label, and let $res$ be an optional Boolean result. If every field of $\sigma$
is at most $B$, then every coordinate of the abstract configuration formed from $l$, the
register-value vector of $\sigma$, and $res$ is at most $B$.
\end{lemma}

\begin{lemma}[Input Pointer Coordinate]
\lean{Complexity.LogProg.UReg.v_0}
\label{Complexity.LogProg.UReg.v_0}
\leanok
\uses{Complexity.LogProg.UReg.v}
\difficulty{1}
For any register state $\sigma$, coordinate $0$ of its register-value vector is its
input pointer $p$.
\end{lemma}

\begin{lemma}[Vertex Count Coordinate]
\lean{Complexity.LogProg.UReg.v_1}
\label{Complexity.LogProg.UReg.v_1}
\leanok
\uses{Complexity.LogProg.UReg.v}
\difficulty{1}
For any register state $\sigma$, coordinate $1$ of its register-value vector is its
vertex count $n$.
\end{lemma}

\begin{lemma}[Outer counter slot]
\lean{Complexity.LogProg.UReg.v_2}
\label{Complexity.LogProg.UReg.v_2}
\leanok
\uses{Complexity.LogProg.UReg.v}
\difficulty{1}
For every register state $\sigma$ and natural number $z$, the value of its register
vector at index $2$ is its outer counter $a$.
\end{lemma}

\begin{lemma}[Inner counter slot]
\lean{Complexity.LogProg.UReg.v_3}
\label{Complexity.LogProg.UReg.v_3}
\leanok
\uses{Complexity.LogProg.UReg.v}
\difficulty{1}
For every register state $\sigma$ and natural number $z$, the value of its register
vector at index $3$ is its inner counter $b$.
\end{lemma}

\begin{lemma}[Current vertex slot]
\lean{Complexity.LogProg.UReg.v_4}
\label{Complexity.LogProg.UReg.v_4}
\leanok
\uses{Complexity.LogProg.UReg.v}
\difficulty{1}
For every register state $\sigma$ and natural number $z$, the value of its register
vector at index $4$ is its current vertex $cur$.
\end{lemma}

\begin{lemma}[Target slot]
\lean{Complexity.LogProg.UReg.v_5}
\label{Complexity.LogProg.UReg.v_5}
\leanok
\uses{Complexity.LogProg.UReg.v}
\difficulty{1}
For every register state $\sigma$ and natural number $z$, the value of its register
vector at index $5$ is its target $t$.
\end{lemma}

\begin{lemma}[Scratch copy slot]
\lean{Complexity.LogProg.UReg.v_6}
\label{Complexity.LogProg.UReg.v_6}
\leanok
\uses{Complexity.LogProg.UReg.v}
\difficulty{1}
For every register state $\sigma$ and natural number $z$, the value of its register
vector at index $6$ is its scratch copy $tmp$.
\end{lemma}

\begin{lemma}[Step count slot]
\lean{Complexity.LogProg.UReg.v_7}
\label{Complexity.LogProg.UReg.v_7}
\leanok
\uses{Complexity.LogProg.UReg.v}
\difficulty{1}
For every register state $\sigma$ and natural number $z$, the value of its register
vector at index $7$ is its step count $st$.
\end{lemma}

\begin{lemma}[Register vector at index eight]
\lean{Complexity.LogProg.UReg.v_8}
\label{Complexity.LogProg.UReg.v_8}
\leanok
\uses{Complexity.LogProg.UReg.v}
\difficulty{1}
For any register state with random-index field $nx$, the eighth entry of its
register-value vector is $nx$.
\end{lemma}

\begin{lemma}[Register vector at index nine]
\lean{Complexity.LogProg.UReg.v_9}
\label{Complexity.LogProg.UReg.v_9}
\leanok
\uses{Complexity.LogProg.UReg.v}
\difficulty{1}
For any register state with saved-field-position field $q$, the ninth entry of its
register-value vector is $q$.
\end{lemma}

\begin{lemma}[Walk Length Coordinate]
\lean{Complexity.LogProg.UReg.v_10}
\label{Complexity.LogProg.UReg.v_10}
\leanok
\uses{Complexity.LogProg.UReg.v}
\difficulty{1}
For any register state $\sigma$, coordinate $10$ of its register-value vector is its
walk length $L$.
\end{lemma}

\begin{lemma}[Candidate Vertex Coordinate]
\lean{Complexity.LogProg.UReg.v_11}
\label{Complexity.LogProg.UReg.v_11}
\leanok
\uses{Complexity.LogProg.UReg.v}
\difficulty{1}
For any register state $\sigma$, coordinate $11$ of its register-value vector is its row
or candidate vertex $u$.
\end{lemma}

\begin{lemma}[Column Coordinate]
\lean{Complexity.LogProg.UReg.v_12}
\label{Complexity.LogProg.UReg.v_12}
\leanok
\uses{Complexity.LogProg.UReg.v}
\difficulty{1}
For any register state $\sigma$, coordinate $12$ of its register-value vector is its
column $w$.
\end{lemma}

\begin{lemma}[Scratch Sum Coordinate]
\lean{Complexity.LogProg.UReg.v_13}
\label{Complexity.LogProg.UReg.v_13}
\leanok
\uses{Complexity.LogProg.UReg.v}
\difficulty{1}
For any register state $\sigma$, coordinate $13$ of its register-value vector is its
scratch sum $e$.
\end{lemma}

\begin{lemma}[Scratch product slot]
\lean{Complexity.LogProg.UReg.v_14}
\label{Complexity.LogProg.UReg.v_14}
\leanok
\uses{Complexity.LogProg.UReg.v}
\difficulty{1}
For every register state $\sigma$ and natural number $z$, the value of its register
vector at index $14$ is its scratch product $f$.
\end{lemma}

\begin{lemma}[Updating the input pointer]
\lean{Complexity.LogProg.UReg.u_0}
\label{Complexity.LogProg.UReg.u_0}
\leanok
\uses{Complexity.LogProg.UReg.v}
\difficulty{4}
For any register state $\sigma$ and natural number $z$, replacing coordinate $0$ of its
register-value vector by $z$ yields the register-value vector of the state obtained by
replacing its input pointer by $z$.
\end{lemma}

\begin{lemma}[Updating the vertex count]
\lean{Complexity.LogProg.UReg.u_1}
\label{Complexity.LogProg.UReg.u_1}
\leanok
\uses{Complexity.LogProg.UReg.v}
\difficulty{4}
For any register state $\sigma$ and natural number $z$, replacing coordinate $1$ of its
register-value vector by $z$ yields the register-value vector of the state obtained by
replacing its vertex count by $z$.
\end{lemma}

\begin{lemma}[Updating the outer counter]
\lean{Complexity.LogProg.UReg.u_2}
\label{Complexity.LogProg.UReg.u_2}
\leanok
\uses{Complexity.LogProg.UReg.v}
\difficulty{4}
For any register state and any natural number $z$, replacing its outer-counter field $a$
by $z$ gives a register-value vector equal to the vector obtained by replacing
coordinate $2$ of the original vector by $z$.
\end{lemma}

\begin{lemma}[Updating the inner counter]
\lean{Complexity.LogProg.UReg.u_3}
\label{Complexity.LogProg.UReg.u_3}
\leanok
\uses{Complexity.LogProg.UReg.v}
\difficulty{4}
For any register state and any natural number $z$, replacing its inner-counter field $b$
by $z$ gives a register-value vector equal to the vector obtained by replacing
coordinate $3$ of the original vector by $z$.
\end{lemma}

\begin{lemma}[Updating the current vertex]
\lean{Complexity.LogProg.UReg.u_4}
\label{Complexity.LogProg.UReg.u_4}
\leanok
\uses{Complexity.LogProg.UReg.v}
\difficulty{4}
For any register state and any natural number $z$, replacing its current-vertex field by
$z$ gives a register-value vector equal to the vector obtained by replacing coordinate
$4$ of the original vector by $z$.
\end{lemma}

\begin{lemma}[Updating the target]
\lean{Complexity.LogProg.UReg.u_5}
\label{Complexity.LogProg.UReg.u_5}
\leanok
\uses{Complexity.LogProg.UReg.v}
\difficulty{4}
For any register state and any natural number $z$, replacing its target field $t$ by $z$
gives a register-value vector equal to the vector obtained by replacing coordinate $5$
of the original vector by $z$.
\end{lemma}

\begin{lemma}[Updating the scratch copy]
\lean{Complexity.LogProg.UReg.u_6}
\label{Complexity.LogProg.UReg.u_6}
\leanok
\uses{Complexity.LogProg.UReg.v}
\difficulty{4}
For any register state and any natural number $z$, replacing its scratch-copy field
$tmp$ by $z$ gives a register-value vector equal to the vector obtained by replacing
coordinate $6$ of the original vector by $z$.
\end{lemma}

\begin{lemma}[Updating the step count]
\lean{Complexity.LogProg.UReg.u_7}
\label{Complexity.LogProg.UReg.u_7}
\leanok
\uses{Complexity.LogProg.UReg.v}
\difficulty{4}
For any register state and any natural number $z$, replacing its step-count field $st$
by $z$ gives a register-value vector equal to the vector obtained by replacing
coordinate $7$ of the original vector by $z$.
\end{lemma}

\begin{lemma}[Updating the random index]
\lean{Complexity.LogProg.UReg.u_8}
\label{Complexity.LogProg.UReg.u_8}
\leanok
\uses{Complexity.LogProg.UReg.v}
\difficulty{4}
For any register state and any natural number $z$, replacing its random-index field $nx$
by $z$ gives a register-value vector equal to the vector obtained by replacing
coordinate $8$ of the original vector by $z$.
\end{lemma}

\begin{lemma}[Updating the Saved Field]
\lean{Complexity.LogProg.UReg.u_9}
\label{Complexity.LogProg.UReg.u_9}
\leanok
\uses{Complexity.LogProg.UReg.v}
\difficulty{4}
For a register state $\sigma$ and a natural number $z$, replacing its saved field
position $q$ by $z$ and then taking its register-value vector gives the same vector as
updating coordinate $9$ of the original vector to $z$.
\end{lemma}

\begin{lemma}[Updating the walk length]
\lean{Complexity.LogProg.UReg.u_10}
\label{Complexity.LogProg.UReg.u_10}
\leanok
\uses{Complexity.LogProg.UReg.v}
\difficulty{4}
For any register state $\sigma$ and natural number $z$, replacing coordinate $10$ of its
register-value vector by $z$ yields the register-value vector of the state obtained by
replacing its walk length by $z$.
\end{lemma}

\begin{lemma}[Updating the candidate vertex]
\lean{Complexity.LogProg.UReg.u_11}
\label{Complexity.LogProg.UReg.u_11}
\leanok
\uses{Complexity.LogProg.UReg.v}
\difficulty{4}
For any register state $\sigma$ and natural number $z$, replacing coordinate $11$ of its
register-value vector by $z$ yields the register-value vector of the state obtained by
replacing its candidate vertex by $z$.
\end{lemma}

\begin{lemma}[Updating the column register]
\lean{Complexity.LogProg.UReg.u_12}
\label{Complexity.LogProg.UReg.u_12}
\leanok
\uses{Complexity.LogProg.UReg.v}
\difficulty{4}
For any register state $\sigma$ and natural number $z$, replacing coordinate $12$ of its
register-value vector by $z$ yields the register-value vector of the state obtained by
replacing its column register by $z$.
\end{lemma}

\begin{lemma}[Updating the scratch sum]
\lean{Complexity.LogProg.UReg.u_13}
\label{Complexity.LogProg.UReg.u_13}
\leanok
\uses{Complexity.LogProg.UReg.v}
\difficulty{4}
For any register state $\sigma$ and natural number $z$, replacing coordinate $13$ of its
register-value vector by $z$ yields the register-value vector of the state obtained by
replacing its scratch sum by $z$.
\end{lemma}

\begin{lemma}[Updating the scratch product]
\lean{Complexity.LogProg.UReg.u_14}
\label{Complexity.LogProg.UReg.u_14}
\leanok
\uses{Complexity.LogProg.UReg.v}
\difficulty{4}
For any register state $\sigma$ and natural number $z$, replacing coordinate $14$ of its
register-value vector by $z$ yields the register-value vector of the state obtained by
replacing its scratch product by $z$.
\end{lemma}

\begin{definition}[Semantics of the path program]
\lean{Complexity.LogProg.UPS}
\statementsource{arora-barak-ch07-randomized-computation}{pdf-26d95080dbde-p016-b002}
\label{Complexity.LogProg.UPS}
\leanok
\uses{Complexity.LogProg.PSem, Complexity.LogProg.RegLe, Complexity.LogProg.ULabel, Complexity.LogProg.UReg, Complexity.LogProg.UReg.v, Complexity.LogProg.pathOr, Complexity.LogProg.upA}
For a Boolean input $x$, a register bound $B$, a control label $\ell$, a register state
$\sigma$, and a rational number $p$, the predicate holds exactly when the specified
program has a path-or execution from initial configuration $(\ell,\sigma)$ on $x$, with
register values bounded by $B$, and probability $p$.
\end{definition}

\begin{lemma}[Invariance under Equal Register States]
\lean{Complexity.LogProg.UPS.of_eq}
\label{Complexity.LogProg.UPS.of_eq}
\leanok
\uses{Complexity.LogProg.ULabel, Complexity.LogProg.UPS, Complexity.LogProg.UReg}
\difficulty{1}
Let $x$ be a Boolean word, let $B$ be a natural number, let $l$ be a program label, let
$p$ be a rational number, and let $\sigma,\sigma'$ be register states. If
$\sigma=\sigma'$ and the probabilistic semantics property with parameter $p$ holds for
the initial configuration formed from $\sigma'$ at label $l$, then it holds for the
initial configuration formed from $\sigma$ at label $l$.
\end{lemma}
